% Literature synthesis v0.2 (2026-10-24), 30 references.
% Include using % Literature synthesis v0.2 (2026-10-24), 30 references.
% Include using % Literature synthesis v0.2 (2026-10-24), 30 references.
% Include using % Literature synthesis v0.2 (2026-10-24), 30 references.
% Include using \input{paper/related_work}; requires booktabs and tabularx.
\section{Related Work}\label{sec:related_work}

\subsection{Threat Model, Dataset Labels, and Email Phishing}
Supervised phishing-email filtering initially emphasized message content, URL, and structural features, as illustrated by Fette \emph{et al.}~\cite{fette2007learning}. Contemporary surveys synthesize machine-learning and NLP approaches~\cite{kavya2025review} and the dual offensive and defensive role of LLMs~\cite{sivaneswaran2026review}. An important evaluation distinction is frequently obscured by aggregate accuracy: predicting phishing versus legitimate intent is not the same as attributing an email to an AI rather than a human author. Likewise, spam cannot automatically be treated as phishing or as legitimate mail. The intended decision, inputs, and class semantics must therefore be fixed before performance comparisons.

Heiding \emph{et al.}~\cite{heiding2024devising} assessed LLM-assisted phishing generation and detection in a human-participant study, while Weinz \emph{et al.}~\cite{weinz2025impact} evaluated emerging phishing and quishing threats in organizational settings. These results motivate realistic security evaluation but do not directly determine the generalization error of a binary text classifier under controlled rewriting.

\subsection{LLM Classification, Behavioral Evidence, and Style}
LLM-based phishing detection already spans multiple inference designs. Zhang \emph{et al.}~\cite{zhang2025sme} evaluated base LLMs on human- and AI-written phishing and legitimate emails for small and midsize enterprises. ChatSpamDetector~\cite{koide2026chatspamdetector} uses message headers and body text to produce phishing assessments and rationales. PRIME~\cite{tian2025prime} aggregates LLM-evaluated social-engineering cues through quantitative and fuzzy scoring, although its reported phishing-only evaluation cannot independently establish benign false-positive rates. Pan \emph{et al.}~\cite{pan2026psychology} quantified persuasion and sensitive-request cues and studied systematic LLM prompting on a large email corpus. MultiPhishGuard~\cite{xue2025multiphishguard} introduced a multi-agent approach combining text, URL, and metadata analysis with adversarial augmentation. Thus, security cues, rationales, adversarially generated email variants, and LLM reasoning are already investigated; any lightweight model must be evaluated against relevant alternatives with comparable access to metadata.

Stylometry represents a related but different target. Opara \emph{et al.}~\cite{opara2025stylometry} investigated commercial spam filters and stylometric cues on AI-generated phishing. CuePhish~\cite{commey2026cuephish} studies interpretable linguistic and stylistic features for distinguishing LLM-generated from non-AI phishing, including adaptive evasion. MAGE~\cite{li2024mage} studies machine-text attribution across generators and domains outside phishing. These studies caution against inferring malicious intent simply from writing style or AI provenance.

\subsection{Rewriting, Generator Transfer, and Corpus Shift}
Ferrell \emph{et al.}~\cite{ferrell2025rephrasing} evaluated detectors after multiple LLMs rephrased phishing and safe messages, reporting degradation and benefits from retraining. Nance \emph{et al.}~\cite{nance2025augmentation} separately explored LLM-generated training augmentation, and Kulal \emph{et al.}~\cite{kulal2025robust} investigated preprocessing against conventional, LLM-generated, and adversarial phishing. Accordingly, ordinary paraphrase testing and augmentation are prior art.

Gutierrez \emph{et al.}~\cite{gutierrez2026crossmodel} studied cross-generator stylometric detection using human and LLM phishing-only samples from three generators. Their cross-model F1 degradation was reduced by recalibration on labeled examples from the target generator. This is important prior evidence for domain-sensitive thresholds, but its positive and negative classes both concern phishing origin. It does not measure benign-AI false alarms, and target-labeled calibration cannot be described as zero-shot generalization.

E-PhishGEN~\cite{pajola2025ephishgen} introduced a multilingual synthetic phishing/legitimate email benchmark to expose weaknesses of older corpora. Bhuiyan and Bhuiyan~\cite{bhuiyan2026fragility} investigated leave-one-corpus-out transfer, prevalence shift, and source artifacts. Their findings connect to the broader shortcut-learning analysis of Geirhos \emph{et al.}~\cite{geirhos2020shortcut} and worst-group robustness methods such as group DRO~\cite{sagawa2020groupdro}. Cross-corpus testing, group robustness, and source-bias diagnosis are consequently necessary controls, not new contributions in themselves.

PhishFuzzer~\cite{toth2026phishfuzzer} provides 3,300 seed messages and 19,800 linked variants covering phishing, spam, and valid email, including sender, URL, and attachment metadata. Parent linkage permits paired analyses, but modifying security-relevant fields can change the correct label. Therefore, a derived message cannot be assumed meaning-preserving merely because its parent ID is known. All derivatives must be grouped before splitting, and source dataset provenance, rather than the superficial generator label, must govern leave-one-generator and joint source-generator tests.

\subsection{Counterfactual Behavioral Testing and Security-Sensitive Edits}
The general distinction between prediction invariance and response to meaningful edits has already been established. CheckList~\cite{ribeiro2020checklist} explicitly defines invariance and directional-expectation tests. Gardner \emph{et al.}~\cite{gardner2020contrast} introduced manually constructed contrast sets to probe local decision boundaries after label-changing modifications. Kaushik \emph{et al.}~\cite{kaushik2020cad} used human-authored, minimally changed counterfactuals to mitigate spurious NLP correlations. Accordingly, bidirectional testing, counterfactual augmentation, and paired sensitivity cannot be claimed as original concepts.

Our proposed phishing-specific investigation must instead determine whether independently verified security-equivalent rewrites and adjudicated security-changing edits can expose or reduce reliance on irrelevant surface features. A changed brand, URL string, or sender field does not by itself establish a legitimate-to-phishing label transition: original context, authority, and actual target must be verified. The experimental question is whether a defined security-semantic annotation protocol and corresponding training strategy outperform equal-budget augmentation, classification-only, and same-label consistency controls on genuinely held-out generators and corpora. No positive finding is assumed.

\subsection{Calibration, Abstention, and Deployment Constraints}
Discrimination and confidence are distinct: Guo \emph{et al.}~\cite{guo2017calibration} analyzed neural calibration, while SelectiveNet~\cite{geifman2019selectivenet} studied rejection under a risk-coverage objective. Tibshirani \emph{et al.}~\cite{tibshirani2019conformal} established conformal methods under specified covariate-shift assumptions, rather than arbitrary distribution-free protection after any unknown shift. Evaluation under rare phishing prevalence needs precision--recall reporting~\cite{saito2015precision}, realized false-positive rates, phishing recall at thresholds frozen from development data, and analyst-review workload.

Arharbi and El Qadi~\cite{arharbi2026shifts} already studied GenAI-induced drift, stability, interpretability, and SHAP/entropy-based gating with lightweight phishing detection. Risk-aware gating alone is therefore prior work. Any additional benefit must be shown under a realistic false-alarm budget without selecting thresholds or rejection policies from the final unseen-generator test labels.

\subsection{Synthesis and Positioning}
Table~\ref{tab:related_comparison} compares the closest work by actual prediction target and experimental scope, rather than pooling accuracy values from incompatible corpora. The literature covers LLM-written phishing, authorship, rephrasing, synthetic benchmarks, stylistic cues, source artifacts, invariant tests, label-changing counterfactuals, and risk-aware decisions. The narrower proposed question is whether a human-adjudicated security-equivalence and security-change protocol supports a measurable improvement on phishing-versus-legitimate classification under parent-disjoint, generator-aware, source-disjoint tests at an explicit false-alarm/review budget. Its novelty and performance remain hypotheses pending dataset and experimental validation.

Table~\ref{tab:our_positioning} makes the proposed contribution's
scope and the exact evidence required for each claim explicit. The
comparison is a research design, not an assertion that earlier studies
never considered these questions. A row will only support a contribution
claim after the stated validation is executed.

\begin{table}[!htbp]
\centering
\caption{Closest related work and exact boundaries of the present research. Statements characterize the cited experiments, not definitive absence claims.}
\label{tab:related_comparison}
\footnotesize
\renewcommand{\arraystretch}{1.12}
\begin{tabularx}{\linewidth}{@{}p{2.65cm}p{2.8cm}X X@{}}
\toprule
Study & Prediction target & Prior capability & Boundary for this study\\
\midrule
Zhang \emph{et al.}~\cite{zhang2025sme} & Phishing/legitimate & Human and AI email, base LLMs & Four origin--intent strata already studied.\\
Koide \emph{et al.}~\cite{koide2026chatspamdetector} & Phishing/legitimate & LLM header/body reasoning & Explainable phishing decisions are prior art.\\
Opara; CuePhish~\cite{opara2025stylometry,commey2026cuephish} & AI vs.\ human phishing signals & Style and adaptive evasion & Authorship differs from phishing intent.\\
Ferrell \emph{et al.}~\cite{ferrell2025rephrasing} & Phishing/legitimate & LLM rewriting and retraining & Rewriting robustness already studied.\\
E-PhishGEN~\cite{pajola2025ephishgen} & Email benchmarking & Multilingual synthetic data & Synthetic data and transfer already studied.\\
Gutierrez \emph{et al.}~\cite{gutierrez2026crossmodel} & Authorship within phishing & Cross-generator and target calibration & Cannot measure legitimate AI-mail FPR.\\
PhishFuzzer~\cite{toth2026phishfuzzer} & Phishing/spam/valid & Parent-linked variants and metadata & Parentage is not semantic equivalence.\\
Bhuiyan and Bhuiyan~\cite{bhuiyan2026fragility} & Phishing/legitimate & Source, prevalence, operational shift & Corpus-shift diagnostics are prior art.\\
CheckList and contrast sets~\cite{ribeiro2020checklist,gardner2020contrast} & General NLP behavioral tests & Invariance and meaningful-change tests & Bidirectional behavior testing is established.\\
Arharbi and El Qadi~\cite{arharbi2026shifts} & Phishing under GenAI shift & Stability and explainable gating & Uncertainty/review mechanisms are prior work.\\
\textbf{Our planned study} & \textbf{Phishing/legitimate} &
\textbf{Adjudicated security-equivalent/changed pairs} &
\textbf{Contribution subject to external validation.}\\
\bottomrule
\end{tabularx}
\end{table}


% PLOS ONE positioning table: research hypotheses, not demonstrated novelty.
\begin{table}[!htbp]
\centering
\caption{Positioning of the proposed research relative to established
approaches. All proposed tests are prospective; meeting an evaluation
criterion has not yet been demonstrated.}
\label{tab:our_positioning}
\footnotesize
\renewcommand{\arraystretch}{1.17}
\begin{tabularx}{\linewidth}{@{}p{2.45cm}X X X@{}}
\toprule
Research strand (prior art) &
Established basis &
Our specific research question / protocol &
Evidence required before a claim \\
\midrule
Authorship and stylistic cues
\cite{opara2025stylometry,gutierrez2026crossmodel,commey2026cuephish} &
AI--human attribution and cross-generator style transfer &
Predict \emph{phishing vs.\ legitimate}; stratify by verified human/AI origin, including AI-written legitimate messages &
Verified labels and provenance; separate benign-AI false-positive rates \\
\addlinespace[2pt]
LLM rewriting and fuzzing
\cite{ferrell2025rephrasing,nance2025augmentation,toth2026phishfuzzer} &
Rephrase-based stress tests, augmentation and parent-linked variants &
Distinguish human-verified, security-equivalent rewrites from adjudicated security-changing edits &
Parent-disjoint pairs; blinded pair-validity audit, disagreements and flip/sensitivity metrics \\
\addlinespace[2pt]
Corpus and generator shift
\cite{pajola2025ephishgen,bhuiyan2026fragility} &
Synthetic/cross-corpus testing and source-artifact diagnoses &
Evaluate separately held-out generator and source, then a joint holdout \emph{only if} data coverage allows &
Exact source/model IDs, zero parent overlap, fixed split and sufficient strata \\
\addlinespace[2pt]
Behavioral counterfactuals
\cite{ribeiro2020checklist,gardner2020contrast,kaushik2020cad} &
Invariance, directional expectation and label-changing revisions &
Test whether a security-specific dual objective adds value over established controls &
Matched-budget augmentation, same-label-only, sensitivity-only and full-model ablations \\
\addlinespace[2pt]
Calibration and risk
\cite{arharbi2026shifts,geifman2019selectivenet} &
Shift monitoring, reject-option and risk--coverage evaluation &
Assess recall at development-chosen false-alarm budgets with optional human review &
Untouched test set, realized FPR, recall, review load and uncertainty intervals \\
\midrule
\textbf{Combined positioning (planned)} &
\textbf{No first-of-its-kind claim} &
\textbf{Security-validated pairs + shift-aware evaluation + risk accounting} &
\textbf{Independent replication and statistically supported benefit or an honest null result} \\
\bottomrule
\end{tabularx}
\end{table}
; requires booktabs and tabularx.
\section{Related Work}\label{sec:related_work}

\subsection{Threat Model, Dataset Labels, and Email Phishing}
Supervised phishing-email filtering initially emphasized message content, URL, and structural features, as illustrated by Fette \emph{et al.}~\cite{fette2007learning}. Contemporary surveys synthesize machine-learning and NLP approaches~\cite{kavya2025review} and the dual offensive and defensive role of LLMs~\cite{sivaneswaran2026review}. An important evaluation distinction is frequently obscured by aggregate accuracy: predicting phishing versus legitimate intent is not the same as attributing an email to an AI rather than a human author. Likewise, spam cannot automatically be treated as phishing or as legitimate mail. The intended decision, inputs, and class semantics must therefore be fixed before performance comparisons.

Heiding \emph{et al.}~\cite{heiding2024devising} assessed LLM-assisted phishing generation and detection in a human-participant study, while Weinz \emph{et al.}~\cite{weinz2025impact} evaluated emerging phishing and quishing threats in organizational settings. These results motivate realistic security evaluation but do not directly determine the generalization error of a binary text classifier under controlled rewriting.

\subsection{LLM Classification, Behavioral Evidence, and Style}
LLM-based phishing detection already spans multiple inference designs. Zhang \emph{et al.}~\cite{zhang2025sme} evaluated base LLMs on human- and AI-written phishing and legitimate emails for small and midsize enterprises. ChatSpamDetector~\cite{koide2026chatspamdetector} uses message headers and body text to produce phishing assessments and rationales. PRIME~\cite{tian2025prime} aggregates LLM-evaluated social-engineering cues through quantitative and fuzzy scoring, although its reported phishing-only evaluation cannot independently establish benign false-positive rates. Pan \emph{et al.}~\cite{pan2026psychology} quantified persuasion and sensitive-request cues and studied systematic LLM prompting on a large email corpus. MultiPhishGuard~\cite{xue2025multiphishguard} introduced a multi-agent approach combining text, URL, and metadata analysis with adversarial augmentation. Thus, security cues, rationales, adversarially generated email variants, and LLM reasoning are already investigated; any lightweight model must be evaluated against relevant alternatives with comparable access to metadata.

Stylometry represents a related but different target. Opara \emph{et al.}~\cite{opara2025stylometry} investigated commercial spam filters and stylometric cues on AI-generated phishing. CuePhish~\cite{commey2026cuephish} studies interpretable linguistic and stylistic features for distinguishing LLM-generated from non-AI phishing, including adaptive evasion. MAGE~\cite{li2024mage} studies machine-text attribution across generators and domains outside phishing. These studies caution against inferring malicious intent simply from writing style or AI provenance.

\subsection{Rewriting, Generator Transfer, and Corpus Shift}
Ferrell \emph{et al.}~\cite{ferrell2025rephrasing} evaluated detectors after multiple LLMs rephrased phishing and safe messages, reporting degradation and benefits from retraining. Nance \emph{et al.}~\cite{nance2025augmentation} separately explored LLM-generated training augmentation, and Kulal \emph{et al.}~\cite{kulal2025robust} investigated preprocessing against conventional, LLM-generated, and adversarial phishing. Accordingly, ordinary paraphrase testing and augmentation are prior art.

Gutierrez \emph{et al.}~\cite{gutierrez2026crossmodel} studied cross-generator stylometric detection using human and LLM phishing-only samples from three generators. Their cross-model F1 degradation was reduced by recalibration on labeled examples from the target generator. This is important prior evidence for domain-sensitive thresholds, but its positive and negative classes both concern phishing origin. It does not measure benign-AI false alarms, and target-labeled calibration cannot be described as zero-shot generalization.

E-PhishGEN~\cite{pajola2025ephishgen} introduced a multilingual synthetic phishing/legitimate email benchmark to expose weaknesses of older corpora. Bhuiyan and Bhuiyan~\cite{bhuiyan2026fragility} investigated leave-one-corpus-out transfer, prevalence shift, and source artifacts. Their findings connect to the broader shortcut-learning analysis of Geirhos \emph{et al.}~\cite{geirhos2020shortcut} and worst-group robustness methods such as group DRO~\cite{sagawa2020groupdro}. Cross-corpus testing, group robustness, and source-bias diagnosis are consequently necessary controls, not new contributions in themselves.

PhishFuzzer~\cite{toth2026phishfuzzer} provides 3,300 seed messages and 19,800 linked variants covering phishing, spam, and valid email, including sender, URL, and attachment metadata. Parent linkage permits paired analyses, but modifying security-relevant fields can change the correct label. Therefore, a derived message cannot be assumed meaning-preserving merely because its parent ID is known. All derivatives must be grouped before splitting, and source dataset provenance, rather than the superficial generator label, must govern leave-one-generator and joint source-generator tests.

\subsection{Counterfactual Behavioral Testing and Security-Sensitive Edits}
The general distinction between prediction invariance and response to meaningful edits has already been established. CheckList~\cite{ribeiro2020checklist} explicitly defines invariance and directional-expectation tests. Gardner \emph{et al.}~\cite{gardner2020contrast} introduced manually constructed contrast sets to probe local decision boundaries after label-changing modifications. Kaushik \emph{et al.}~\cite{kaushik2020cad} used human-authored, minimally changed counterfactuals to mitigate spurious NLP correlations. Accordingly, bidirectional testing, counterfactual augmentation, and paired sensitivity cannot be claimed as original concepts.

Our proposed phishing-specific investigation must instead determine whether independently verified security-equivalent rewrites and adjudicated security-changing edits can expose or reduce reliance on irrelevant surface features. A changed brand, URL string, or sender field does not by itself establish a legitimate-to-phishing label transition: original context, authority, and actual target must be verified. The experimental question is whether a defined security-semantic annotation protocol and corresponding training strategy outperform equal-budget augmentation, classification-only, and same-label consistency controls on genuinely held-out generators and corpora. No positive finding is assumed.

\subsection{Calibration, Abstention, and Deployment Constraints}
Discrimination and confidence are distinct: Guo \emph{et al.}~\cite{guo2017calibration} analyzed neural calibration, while SelectiveNet~\cite{geifman2019selectivenet} studied rejection under a risk-coverage objective. Tibshirani \emph{et al.}~\cite{tibshirani2019conformal} established conformal methods under specified covariate-shift assumptions, rather than arbitrary distribution-free protection after any unknown shift. Evaluation under rare phishing prevalence needs precision--recall reporting~\cite{saito2015precision}, realized false-positive rates, phishing recall at thresholds frozen from development data, and analyst-review workload.

Arharbi and El Qadi~\cite{arharbi2026shifts} already studied GenAI-induced drift, stability, interpretability, and SHAP/entropy-based gating with lightweight phishing detection. Risk-aware gating alone is therefore prior work. Any additional benefit must be shown under a realistic false-alarm budget without selecting thresholds or rejection policies from the final unseen-generator test labels.

\subsection{Synthesis and Positioning}
Table~\ref{tab:related_comparison} compares the closest work by actual prediction target and experimental scope, rather than pooling accuracy values from incompatible corpora. The literature covers LLM-written phishing, authorship, rephrasing, synthetic benchmarks, stylistic cues, source artifacts, invariant tests, label-changing counterfactuals, and risk-aware decisions. The narrower proposed question is whether a human-adjudicated security-equivalence and security-change protocol supports a measurable improvement on phishing-versus-legitimate classification under parent-disjoint, generator-aware, source-disjoint tests at an explicit false-alarm/review budget. Its novelty and performance remain hypotheses pending dataset and experimental validation.

Table~\ref{tab:our_positioning} makes the proposed contribution's
scope and the exact evidence required for each claim explicit. The
comparison is a research design, not an assertion that earlier studies
never considered these questions. A row will only support a contribution
claim after the stated validation is executed.

\begin{table}[!htbp]
\centering
\caption{Closest related work and exact boundaries of the present research. Statements characterize the cited experiments, not definitive absence claims.}
\label{tab:related_comparison}
\footnotesize
\renewcommand{\arraystretch}{1.12}
\begin{tabularx}{\linewidth}{@{}p{2.65cm}p{2.8cm}X X@{}}
\toprule
Study & Prediction target & Prior capability & Boundary for this study\\
\midrule
Zhang \emph{et al.}~\cite{zhang2025sme} & Phishing/legitimate & Human and AI email, base LLMs & Four origin--intent strata already studied.\\
Koide \emph{et al.}~\cite{koide2026chatspamdetector} & Phishing/legitimate & LLM header/body reasoning & Explainable phishing decisions are prior art.\\
Opara; CuePhish~\cite{opara2025stylometry,commey2026cuephish} & AI vs.\ human phishing signals & Style and adaptive evasion & Authorship differs from phishing intent.\\
Ferrell \emph{et al.}~\cite{ferrell2025rephrasing} & Phishing/legitimate & LLM rewriting and retraining & Rewriting robustness already studied.\\
E-PhishGEN~\cite{pajola2025ephishgen} & Email benchmarking & Multilingual synthetic data & Synthetic data and transfer already studied.\\
Gutierrez \emph{et al.}~\cite{gutierrez2026crossmodel} & Authorship within phishing & Cross-generator and target calibration & Cannot measure legitimate AI-mail FPR.\\
PhishFuzzer~\cite{toth2026phishfuzzer} & Phishing/spam/valid & Parent-linked variants and metadata & Parentage is not semantic equivalence.\\
Bhuiyan and Bhuiyan~\cite{bhuiyan2026fragility} & Phishing/legitimate & Source, prevalence, operational shift & Corpus-shift diagnostics are prior art.\\
CheckList and contrast sets~\cite{ribeiro2020checklist,gardner2020contrast} & General NLP behavioral tests & Invariance and meaningful-change tests & Bidirectional behavior testing is established.\\
Arharbi and El Qadi~\cite{arharbi2026shifts} & Phishing under GenAI shift & Stability and explainable gating & Uncertainty/review mechanisms are prior work.\\
\textbf{Our planned study} & \textbf{Phishing/legitimate} &
\textbf{Adjudicated security-equivalent/changed pairs} &
\textbf{Contribution subject to external validation.}\\
\bottomrule
\end{tabularx}
\end{table}


% PLOS ONE positioning table: research hypotheses, not demonstrated novelty.
\begin{table}[!htbp]
\centering
\caption{Positioning of the proposed research relative to established
approaches. All proposed tests are prospective; meeting an evaluation
criterion has not yet been demonstrated.}
\label{tab:our_positioning}
\footnotesize
\renewcommand{\arraystretch}{1.17}
\begin{tabularx}{\linewidth}{@{}p{2.45cm}X X X@{}}
\toprule
Research strand (prior art) &
Established basis &
Our specific research question / protocol &
Evidence required before a claim \\
\midrule
Authorship and stylistic cues
\cite{opara2025stylometry,gutierrez2026crossmodel,commey2026cuephish} &
AI--human attribution and cross-generator style transfer &
Predict \emph{phishing vs.\ legitimate}; stratify by verified human/AI origin, including AI-written legitimate messages &
Verified labels and provenance; separate benign-AI false-positive rates \\
\addlinespace[2pt]
LLM rewriting and fuzzing
\cite{ferrell2025rephrasing,nance2025augmentation,toth2026phishfuzzer} &
Rephrase-based stress tests, augmentation and parent-linked variants &
Distinguish human-verified, security-equivalent rewrites from adjudicated security-changing edits &
Parent-disjoint pairs; blinded pair-validity audit, disagreements and flip/sensitivity metrics \\
\addlinespace[2pt]
Corpus and generator shift
\cite{pajola2025ephishgen,bhuiyan2026fragility} &
Synthetic/cross-corpus testing and source-artifact diagnoses &
Evaluate separately held-out generator and source, then a joint holdout \emph{only if} data coverage allows &
Exact source/model IDs, zero parent overlap, fixed split and sufficient strata \\
\addlinespace[2pt]
Behavioral counterfactuals
\cite{ribeiro2020checklist,gardner2020contrast,kaushik2020cad} &
Invariance, directional expectation and label-changing revisions &
Test whether a security-specific dual objective adds value over established controls &
Matched-budget augmentation, same-label-only, sensitivity-only and full-model ablations \\
\addlinespace[2pt]
Calibration and risk
\cite{arharbi2026shifts,geifman2019selectivenet} &
Shift monitoring, reject-option and risk--coverage evaluation &
Assess recall at development-chosen false-alarm budgets with optional human review &
Untouched test set, realized FPR, recall, review load and uncertainty intervals \\
\midrule
\textbf{Combined positioning (planned)} &
\textbf{No first-of-its-kind claim} &
\textbf{Security-validated pairs + shift-aware evaluation + risk accounting} &
\textbf{Independent replication and statistically supported benefit or an honest null result} \\
\bottomrule
\end{tabularx}
\end{table}
; requires booktabs and tabularx.
\section{Related Work}\label{sec:related_work}

\subsection{Threat Model, Dataset Labels, and Email Phishing}
Supervised phishing-email filtering initially emphasized message content, URL, and structural features, as illustrated by Fette \emph{et al.}~\cite{fette2007learning}. Contemporary surveys synthesize machine-learning and NLP approaches~\cite{kavya2025review} and the dual offensive and defensive role of LLMs~\cite{sivaneswaran2026review}. An important evaluation distinction is frequently obscured by aggregate accuracy: predicting phishing versus legitimate intent is not the same as attributing an email to an AI rather than a human author. Likewise, spam cannot automatically be treated as phishing or as legitimate mail. The intended decision, inputs, and class semantics must therefore be fixed before performance comparisons.

Heiding \emph{et al.}~\cite{heiding2024devising} assessed LLM-assisted phishing generation and detection in a human-participant study, while Weinz \emph{et al.}~\cite{weinz2025impact} evaluated emerging phishing and quishing threats in organizational settings. These results motivate realistic security evaluation but do not directly determine the generalization error of a binary text classifier under controlled rewriting.

\subsection{LLM Classification, Behavioral Evidence, and Style}
LLM-based phishing detection already spans multiple inference designs. Zhang \emph{et al.}~\cite{zhang2025sme} evaluated base LLMs on human- and AI-written phishing and legitimate emails for small and midsize enterprises. ChatSpamDetector~\cite{koide2026chatspamdetector} uses message headers and body text to produce phishing assessments and rationales. PRIME~\cite{tian2025prime} aggregates LLM-evaluated social-engineering cues through quantitative and fuzzy scoring, although its reported phishing-only evaluation cannot independently establish benign false-positive rates. Pan \emph{et al.}~\cite{pan2026psychology} quantified persuasion and sensitive-request cues and studied systematic LLM prompting on a large email corpus. MultiPhishGuard~\cite{xue2025multiphishguard} introduced a multi-agent approach combining text, URL, and metadata analysis with adversarial augmentation. Thus, security cues, rationales, adversarially generated email variants, and LLM reasoning are already investigated; any lightweight model must be evaluated against relevant alternatives with comparable access to metadata.

Stylometry represents a related but different target. Opara \emph{et al.}~\cite{opara2025stylometry} investigated commercial spam filters and stylometric cues on AI-generated phishing. CuePhish~\cite{commey2026cuephish} studies interpretable linguistic and stylistic features for distinguishing LLM-generated from non-AI phishing, including adaptive evasion. MAGE~\cite{li2024mage} studies machine-text attribution across generators and domains outside phishing. These studies caution against inferring malicious intent simply from writing style or AI provenance.

\subsection{Rewriting, Generator Transfer, and Corpus Shift}
Ferrell \emph{et al.}~\cite{ferrell2025rephrasing} evaluated detectors after multiple LLMs rephrased phishing and safe messages, reporting degradation and benefits from retraining. Nance \emph{et al.}~\cite{nance2025augmentation} separately explored LLM-generated training augmentation, and Kulal \emph{et al.}~\cite{kulal2025robust} investigated preprocessing against conventional, LLM-generated, and adversarial phishing. Accordingly, ordinary paraphrase testing and augmentation are prior art.

Gutierrez \emph{et al.}~\cite{gutierrez2026crossmodel} studied cross-generator stylometric detection using human and LLM phishing-only samples from three generators. Their cross-model F1 degradation was reduced by recalibration on labeled examples from the target generator. This is important prior evidence for domain-sensitive thresholds, but its positive and negative classes both concern phishing origin. It does not measure benign-AI false alarms, and target-labeled calibration cannot be described as zero-shot generalization.

E-PhishGEN~\cite{pajola2025ephishgen} introduced a multilingual synthetic phishing/legitimate email benchmark to expose weaknesses of older corpora. Bhuiyan and Bhuiyan~\cite{bhuiyan2026fragility} investigated leave-one-corpus-out transfer, prevalence shift, and source artifacts. Their findings connect to the broader shortcut-learning analysis of Geirhos \emph{et al.}~\cite{geirhos2020shortcut} and worst-group robustness methods such as group DRO~\cite{sagawa2020groupdro}. Cross-corpus testing, group robustness, and source-bias diagnosis are consequently necessary controls, not new contributions in themselves.

PhishFuzzer~\cite{toth2026phishfuzzer} provides 3,300 seed messages and 19,800 linked variants covering phishing, spam, and valid email, including sender, URL, and attachment metadata. Parent linkage permits paired analyses, but modifying security-relevant fields can change the correct label. Therefore, a derived message cannot be assumed meaning-preserving merely because its parent ID is known. All derivatives must be grouped before splitting, and source dataset provenance, rather than the superficial generator label, must govern leave-one-generator and joint source-generator tests.

\subsection{Counterfactual Behavioral Testing and Security-Sensitive Edits}
The general distinction between prediction invariance and response to meaningful edits has already been established. CheckList~\cite{ribeiro2020checklist} explicitly defines invariance and directional-expectation tests. Gardner \emph{et al.}~\cite{gardner2020contrast} introduced manually constructed contrast sets to probe local decision boundaries after label-changing modifications. Kaushik \emph{et al.}~\cite{kaushik2020cad} used human-authored, minimally changed counterfactuals to mitigate spurious NLP correlations. Accordingly, bidirectional testing, counterfactual augmentation, and paired sensitivity cannot be claimed as original concepts.

Our proposed phishing-specific investigation must instead determine whether independently verified security-equivalent rewrites and adjudicated security-changing edits can expose or reduce reliance on irrelevant surface features. A changed brand, URL string, or sender field does not by itself establish a legitimate-to-phishing label transition: original context, authority, and actual target must be verified. The experimental question is whether a defined security-semantic annotation protocol and corresponding training strategy outperform equal-budget augmentation, classification-only, and same-label consistency controls on genuinely held-out generators and corpora. No positive finding is assumed.

\subsection{Calibration, Abstention, and Deployment Constraints}
Discrimination and confidence are distinct: Guo \emph{et al.}~\cite{guo2017calibration} analyzed neural calibration, while SelectiveNet~\cite{geifman2019selectivenet} studied rejection under a risk-coverage objective. Tibshirani \emph{et al.}~\cite{tibshirani2019conformal} established conformal methods under specified covariate-shift assumptions, rather than arbitrary distribution-free protection after any unknown shift. Evaluation under rare phishing prevalence needs precision--recall reporting~\cite{saito2015precision}, realized false-positive rates, phishing recall at thresholds frozen from development data, and analyst-review workload.

Arharbi and El Qadi~\cite{arharbi2026shifts} already studied GenAI-induced drift, stability, interpretability, and SHAP/entropy-based gating with lightweight phishing detection. Risk-aware gating alone is therefore prior work. Any additional benefit must be shown under a realistic false-alarm budget without selecting thresholds or rejection policies from the final unseen-generator test labels.

\subsection{Synthesis and Positioning}
Table~\ref{tab:related_comparison} compares the closest work by actual prediction target and experimental scope, rather than pooling accuracy values from incompatible corpora. The literature covers LLM-written phishing, authorship, rephrasing, synthetic benchmarks, stylistic cues, source artifacts, invariant tests, label-changing counterfactuals, and risk-aware decisions. The narrower proposed question is whether a human-adjudicated security-equivalence and security-change protocol supports a measurable improvement on phishing-versus-legitimate classification under parent-disjoint, generator-aware, source-disjoint tests at an explicit false-alarm/review budget. Its novelty and performance remain hypotheses pending dataset and experimental validation.

Table~\ref{tab:our_positioning} makes the proposed contribution's
scope and the exact evidence required for each claim explicit. The
comparison is a research design, not an assertion that earlier studies
never considered these questions. A row will only support a contribution
claim after the stated validation is executed.

\begin{table}[!htbp]
\centering
\caption{Closest related work and exact boundaries of the present research. Statements characterize the cited experiments, not definitive absence claims.}
\label{tab:related_comparison}
\footnotesize
\renewcommand{\arraystretch}{1.12}
\begin{tabularx}{\linewidth}{@{}p{2.65cm}p{2.8cm}X X@{}}
\toprule
Study & Prediction target & Prior capability & Boundary for this study\\
\midrule
Zhang \emph{et al.}~\cite{zhang2025sme} & Phishing/legitimate & Human and AI email, base LLMs & Four origin--intent strata already studied.\\
Koide \emph{et al.}~\cite{koide2026chatspamdetector} & Phishing/legitimate & LLM header/body reasoning & Explainable phishing decisions are prior art.\\
Opara; CuePhish~\cite{opara2025stylometry,commey2026cuephish} & AI vs.\ human phishing signals & Style and adaptive evasion & Authorship differs from phishing intent.\\
Ferrell \emph{et al.}~\cite{ferrell2025rephrasing} & Phishing/legitimate & LLM rewriting and retraining & Rewriting robustness already studied.\\
E-PhishGEN~\cite{pajola2025ephishgen} & Email benchmarking & Multilingual synthetic data & Synthetic data and transfer already studied.\\
Gutierrez \emph{et al.}~\cite{gutierrez2026crossmodel} & Authorship within phishing & Cross-generator and target calibration & Cannot measure legitimate AI-mail FPR.\\
PhishFuzzer~\cite{toth2026phishfuzzer} & Phishing/spam/valid & Parent-linked variants and metadata & Parentage is not semantic equivalence.\\
Bhuiyan and Bhuiyan~\cite{bhuiyan2026fragility} & Phishing/legitimate & Source, prevalence, operational shift & Corpus-shift diagnostics are prior art.\\
CheckList and contrast sets~\cite{ribeiro2020checklist,gardner2020contrast} & General NLP behavioral tests & Invariance and meaningful-change tests & Bidirectional behavior testing is established.\\
Arharbi and El Qadi~\cite{arharbi2026shifts} & Phishing under GenAI shift & Stability and explainable gating & Uncertainty/review mechanisms are prior work.\\
\textbf{Our planned study} & \textbf{Phishing/legitimate} &
\textbf{Adjudicated security-equivalent/changed pairs} &
\textbf{Contribution subject to external validation.}\\
\bottomrule
\end{tabularx}
\end{table}


% PLOS ONE positioning table: research hypotheses, not demonstrated novelty.
\begin{table}[!htbp]
\centering
\caption{Positioning of the proposed research relative to established
approaches. All proposed tests are prospective; meeting an evaluation
criterion has not yet been demonstrated.}
\label{tab:our_positioning}
\footnotesize
\renewcommand{\arraystretch}{1.17}
\begin{tabularx}{\linewidth}{@{}p{2.45cm}X X X@{}}
\toprule
Research strand (prior art) &
Established basis &
Our specific research question / protocol &
Evidence required before a claim \\
\midrule
Authorship and stylistic cues
\cite{opara2025stylometry,gutierrez2026crossmodel,commey2026cuephish} &
AI--human attribution and cross-generator style transfer &
Predict \emph{phishing vs.\ legitimate}; stratify by verified human/AI origin, including AI-written legitimate messages &
Verified labels and provenance; separate benign-AI false-positive rates \\
\addlinespace[2pt]
LLM rewriting and fuzzing
\cite{ferrell2025rephrasing,nance2025augmentation,toth2026phishfuzzer} &
Rephrase-based stress tests, augmentation and parent-linked variants &
Distinguish human-verified, security-equivalent rewrites from adjudicated security-changing edits &
Parent-disjoint pairs; blinded pair-validity audit, disagreements and flip/sensitivity metrics \\
\addlinespace[2pt]
Corpus and generator shift
\cite{pajola2025ephishgen,bhuiyan2026fragility} &
Synthetic/cross-corpus testing and source-artifact diagnoses &
Evaluate separately held-out generator and source, then a joint holdout \emph{only if} data coverage allows &
Exact source/model IDs, zero parent overlap, fixed split and sufficient strata \\
\addlinespace[2pt]
Behavioral counterfactuals
\cite{ribeiro2020checklist,gardner2020contrast,kaushik2020cad} &
Invariance, directional expectation and label-changing revisions &
Test whether a security-specific dual objective adds value over established controls &
Matched-budget augmentation, same-label-only, sensitivity-only and full-model ablations \\
\addlinespace[2pt]
Calibration and risk
\cite{arharbi2026shifts,geifman2019selectivenet} &
Shift monitoring, reject-option and risk--coverage evaluation &
Assess recall at development-chosen false-alarm budgets with optional human review &
Untouched test set, realized FPR, recall, review load and uncertainty intervals \\
\midrule
\textbf{Combined positioning (planned)} &
\textbf{No first-of-its-kind claim} &
\textbf{Security-validated pairs + shift-aware evaluation + risk accounting} &
\textbf{Independent replication and statistically supported benefit or an honest null result} \\
\bottomrule
\end{tabularx}
\end{table}
; requires booktabs and tabularx.
\section{Related Work}\label{sec:related_work}

\subsection{Threat Model, Dataset Labels, and Email Phishing}
Supervised phishing-email filtering initially emphasized message content, URL, and structural features, as illustrated by Fette \emph{et al.}~\cite{fette2007learning}. Contemporary surveys synthesize machine-learning and NLP approaches~\cite{kavya2025review} and the dual offensive and defensive role of LLMs~\cite{sivaneswaran2026review}. An important evaluation distinction is frequently obscured by aggregate accuracy: predicting phishing versus legitimate intent is not the same as attributing an email to an AI rather than a human author. Likewise, spam cannot automatically be treated as phishing or as legitimate mail. The intended decision, inputs, and class semantics must therefore be fixed before performance comparisons.

Heiding \emph{et al.}~\cite{heiding2024devising} assessed LLM-assisted phishing generation and detection in a human-participant study, while Weinz \emph{et al.}~\cite{weinz2025impact} evaluated emerging phishing and quishing threats in organizational settings. These results motivate realistic security evaluation but do not directly determine the generalization error of a binary text classifier under controlled rewriting.

\subsection{LLM Classification, Behavioral Evidence, and Style}
LLM-based phishing detection already spans multiple inference designs. Zhang \emph{et al.}~\cite{zhang2025sme} evaluated base LLMs on human- and AI-written phishing and legitimate emails for small and midsize enterprises. ChatSpamDetector~\cite{koide2026chatspamdetector} uses message headers and body text to produce phishing assessments and rationales. PRIME~\cite{tian2025prime} aggregates LLM-evaluated social-engineering cues through quantitative and fuzzy scoring, although its reported phishing-only evaluation cannot independently establish benign false-positive rates. Pan \emph{et al.}~\cite{pan2026psychology} quantified persuasion and sensitive-request cues and studied systematic LLM prompting on a large email corpus. MultiPhishGuard~\cite{xue2025multiphishguard} introduced a multi-agent approach combining text, URL, and metadata analysis with adversarial augmentation. Thus, security cues, rationales, adversarially generated email variants, and LLM reasoning are already investigated; any lightweight model must be evaluated against relevant alternatives with comparable access to metadata.

Stylometry represents a related but different target. Opara \emph{et al.}~\cite{opara2025stylometry} investigated commercial spam filters and stylometric cues on AI-generated phishing. CuePhish~\cite{commey2026cuephish} studies interpretable linguistic and stylistic features for distinguishing LLM-generated from non-AI phishing, including adaptive evasion. MAGE~\cite{li2024mage} studies machine-text attribution across generators and domains outside phishing. These studies caution against inferring malicious intent simply from writing style or AI provenance.

\subsection{Rewriting, Generator Transfer, and Corpus Shift}
Ferrell \emph{et al.}~\cite{ferrell2025rephrasing} evaluated detectors after multiple LLMs rephrased phishing and safe messages, reporting degradation and benefits from retraining. Nance \emph{et al.}~\cite{nance2025augmentation} separately explored LLM-generated training augmentation, and Kulal \emph{et al.}~\cite{kulal2025robust} investigated preprocessing against conventional, LLM-generated, and adversarial phishing. Accordingly, ordinary paraphrase testing and augmentation are prior art.

Gutierrez \emph{et al.}~\cite{gutierrez2026crossmodel} studied cross-generator stylometric detection using human and LLM phishing-only samples from three generators. Their cross-model F1 degradation was reduced by recalibration on labeled examples from the target generator. This is important prior evidence for domain-sensitive thresholds, but its positive and negative classes both concern phishing origin. It does not measure benign-AI false alarms, and target-labeled calibration cannot be described as zero-shot generalization.

E-PhishGEN~\cite{pajola2025ephishgen} introduced a multilingual synthetic phishing/legitimate email benchmark to expose weaknesses of older corpora. Bhuiyan and Bhuiyan~\cite{bhuiyan2026fragility} investigated leave-one-corpus-out transfer, prevalence shift, and source artifacts. Their findings connect to the broader shortcut-learning analysis of Geirhos \emph{et al.}~\cite{geirhos2020shortcut} and worst-group robustness methods such as group DRO~\cite{sagawa2020groupdro}. Cross-corpus testing, group robustness, and source-bias diagnosis are consequently necessary controls, not new contributions in themselves.

PhishFuzzer~\cite{toth2026phishfuzzer} provides 3,300 seed messages and 19,800 linked variants covering phishing, spam, and valid email, including sender, URL, and attachment metadata. Parent linkage permits paired analyses, but modifying security-relevant fields can change the correct label. Therefore, a derived message cannot be assumed meaning-preserving merely because its parent ID is known. All derivatives must be grouped before splitting, and source dataset provenance, rather than the superficial generator label, must govern leave-one-generator and joint source-generator tests.

\subsection{Counterfactual Behavioral Testing and Security-Sensitive Edits}
The general distinction between prediction invariance and response to meaningful edits has already been established. CheckList~\cite{ribeiro2020checklist} explicitly defines invariance and directional-expectation tests. Gardner \emph{et al.}~\cite{gardner2020contrast} introduced manually constructed contrast sets to probe local decision boundaries after label-changing modifications. Kaushik \emph{et al.}~\cite{kaushik2020cad} used human-authored, minimally changed counterfactuals to mitigate spurious NLP correlations. Accordingly, bidirectional testing, counterfactual augmentation, and paired sensitivity cannot be claimed as original concepts.

Our proposed phishing-specific investigation must instead determine whether independently verified security-equivalent rewrites and adjudicated security-changing edits can expose or reduce reliance on irrelevant surface features. A changed brand, URL string, or sender field does not by itself establish a legitimate-to-phishing label transition: original context, authority, and actual target must be verified. The experimental question is whether a defined security-semantic annotation protocol and corresponding training strategy outperform equal-budget augmentation, classification-only, and same-label consistency controls on genuinely held-out generators and corpora. No positive finding is assumed.

\subsection{Calibration, Abstention, and Deployment Constraints}
Discrimination and confidence are distinct: Guo \emph{et al.}~\cite{guo2017calibration} analyzed neural calibration, while SelectiveNet~\cite{geifman2019selectivenet} studied rejection under a risk-coverage objective. Tibshirani \emph{et al.}~\cite{tibshirani2019conformal} established conformal methods under specified covariate-shift assumptions, rather than arbitrary distribution-free protection after any unknown shift. Evaluation under rare phishing prevalence needs precision--recall reporting~\cite{saito2015precision}, realized false-positive rates, phishing recall at thresholds frozen from development data, and analyst-review workload.

Arharbi and El Qadi~\cite{arharbi2026shifts} already studied GenAI-induced drift, stability, interpretability, and SHAP/entropy-based gating with lightweight phishing detection. Risk-aware gating alone is therefore prior work. Any additional benefit must be shown under a realistic false-alarm budget without selecting thresholds or rejection policies from the final unseen-generator test labels.

\subsection{Synthesis and Positioning}
Table~\ref{tab:related_comparison} compares the closest work by actual prediction target and experimental scope, rather than pooling accuracy values from incompatible corpora. The literature covers LLM-written phishing, authorship, rephrasing, synthetic benchmarks, stylistic cues, source artifacts, invariant tests, label-changing counterfactuals, and risk-aware decisions. The narrower proposed question is whether a human-adjudicated security-equivalence and security-change protocol supports a measurable improvement on phishing-versus-legitimate classification under parent-disjoint, generator-aware, source-disjoint tests at an explicit false-alarm/review budget. Its novelty and performance remain hypotheses pending dataset and experimental validation.

Table~\ref{tab:our_positioning} makes the proposed contribution's
scope and the exact evidence required for each claim explicit. The
comparison is a research design, not an assertion that earlier studies
never considered these questions. A row will only support a contribution
claim after the stated validation is executed.

\begin{table}[!htbp]
\centering
\caption{Closest related work and exact boundaries of the present research. Statements characterize the cited experiments, not definitive absence claims.}
\label{tab:related_comparison}
\footnotesize
\renewcommand{\arraystretch}{1.12}
\begin{tabularx}{\linewidth}{@{}p{2.65cm}p{2.8cm}X X@{}}
\toprule
Study & Prediction target & Prior capability & Boundary for this study\\
\midrule
Zhang \emph{et al.}~\cite{zhang2025sme} & Phishing/legitimate & Human and AI email, base LLMs & Four origin--intent strata already studied.\\
Koide \emph{et al.}~\cite{koide2026chatspamdetector} & Phishing/legitimate & LLM header/body reasoning & Explainable phishing decisions are prior art.\\
Opara; CuePhish~\cite{opara2025stylometry,commey2026cuephish} & AI vs.\ human phishing signals & Style and adaptive evasion & Authorship differs from phishing intent.\\
Ferrell \emph{et al.}~\cite{ferrell2025rephrasing} & Phishing/legitimate & LLM rewriting and retraining & Rewriting robustness already studied.\\
E-PhishGEN~\cite{pajola2025ephishgen} & Email benchmarking & Multilingual synthetic data & Synthetic data and transfer already studied.\\
Gutierrez \emph{et al.}~\cite{gutierrez2026crossmodel} & Authorship within phishing & Cross-generator and target calibration & Cannot measure legitimate AI-mail FPR.\\
PhishFuzzer~\cite{toth2026phishfuzzer} & Phishing/spam/valid & Parent-linked variants and metadata & Parentage is not semantic equivalence.\\
Bhuiyan and Bhuiyan~\cite{bhuiyan2026fragility} & Phishing/legitimate & Source, prevalence, operational shift & Corpus-shift diagnostics are prior art.\\
CheckList and contrast sets~\cite{ribeiro2020checklist,gardner2020contrast} & General NLP behavioral tests & Invariance and meaningful-change tests & Bidirectional behavior testing is established.\\
Arharbi and El Qadi~\cite{arharbi2026shifts} & Phishing under GenAI shift & Stability and explainable gating & Uncertainty/review mechanisms are prior work.\\
\textbf{Our planned study} & \textbf{Phishing/legitimate} &
\textbf{Adjudicated security-equivalent/changed pairs} &
\textbf{Contribution subject to external validation.}\\
\bottomrule
\end{tabularx}
\end{table}


% PLOS ONE positioning table: research hypotheses, not demonstrated novelty.
\begin{table}[!htbp]
\centering
\caption{Positioning of the proposed research relative to established
approaches. All proposed tests are prospective; meeting an evaluation
criterion has not yet been demonstrated.}
\label{tab:our_positioning}
\footnotesize
\renewcommand{\arraystretch}{1.17}
\begin{tabularx}{\linewidth}{@{}p{2.45cm}X X X@{}}
\toprule
Research strand (prior art) &
Established basis &
Our specific research question / protocol &
Evidence required before a claim \\
\midrule
Authorship and stylistic cues
\cite{opara2025stylometry,gutierrez2026crossmodel,commey2026cuephish} &
AI--human attribution and cross-generator style transfer &
Predict \emph{phishing vs.\ legitimate}; stratify by verified human/AI origin, including AI-written legitimate messages &
Verified labels and provenance; separate benign-AI false-positive rates \\
\addlinespace[2pt]
LLM rewriting and fuzzing
\cite{ferrell2025rephrasing,nance2025augmentation,toth2026phishfuzzer} &
Rephrase-based stress tests, augmentation and parent-linked variants &
Distinguish human-verified, security-equivalent rewrites from adjudicated security-changing edits &
Parent-disjoint pairs; blinded pair-validity audit, disagreements and flip/sensitivity metrics \\
\addlinespace[2pt]
Corpus and generator shift
\cite{pajola2025ephishgen,bhuiyan2026fragility} &
Synthetic/cross-corpus testing and source-artifact diagnoses &
Evaluate separately held-out generator and source, then a joint holdout \emph{only if} data coverage allows &
Exact source/model IDs, zero parent overlap, fixed split and sufficient strata \\
\addlinespace[2pt]
Behavioral counterfactuals
\cite{ribeiro2020checklist,gardner2020contrast,kaushik2020cad} &
Invariance, directional expectation and label-changing revisions &
Test whether a security-specific dual objective adds value over established controls &
Matched-budget augmentation, same-label-only, sensitivity-only and full-model ablations \\
\addlinespace[2pt]
Calibration and risk
\cite{arharbi2026shifts,geifman2019selectivenet} &
Shift monitoring, reject-option and risk--coverage evaluation &
Assess recall at development-chosen false-alarm budgets with optional human review &
Untouched test set, realized FPR, recall, review load and uncertainty intervals \\
\midrule
\textbf{Combined positioning (planned)} &
\textbf{No first-of-its-kind claim} &
\textbf{Security-validated pairs + shift-aware evaluation + risk accounting} &
\textbf{Independent replication and statistically supported benefit or an honest null result} \\
\bottomrule
\end{tabularx}
\end{table}
